% =============================================================
% Capitulo 10 - Descripcion del modelo experimental
% =============================================================
\chapter{Descripci\'on del modelo experimental}
\label{cap:modelo_experimental}

Este cap\'itulo describe el modelo experimental terminado, integrando
el circuito de detecci\'on de picos (Cap\'itulo~\ref{cap:diseno_deteccion_picos}),
la placa de control (Cap\'itulo~\ref{cap:placa_control}), el firmware
del microcontrolador (Cap\'itulo~\ref{cap:firmware}) y el software HMI
(Cap\'itulo~\ref{cap:software_hmi}) en un \'unico sistema operativo sobre
el HSRL del Observatorio de Pilar. Al momento de redacci\'on de este
informe, la integraci\'on final de la placa de control se encuentra en
curso \cite{informeAbril2026}, por lo que las secciones siguientes se
completar\'an una vez finalizado el ensamble y la puesta a punto del
sistema.

\section{Integraci\'on del sistema completo}
\label{sec:integracion_sistema}

PLACEHOLDER: describir la integraci\'on f\'isica del sistema completo
sobre el HSRL de Pilar, en reemplazo del conjunto osciloscopio-PC
descrito en la Secci\'on~\ref{sec:antecedentes}, una vez finalizada la
placa de control.

\section{Manual de operaci\'on}
\label{sec:manual_operacion}

PLACEHOLDER: incorporar el manual de operaci\'on del sistema (conexionado,
encendido, selecci\'on de modo de operaci\'on seg\'un lo descrito en la
Secci\'on~\ref{sec:modos_operacion}, uso del HMI) una vez validado el
modelo experimental completo.

\section{Procedimiento de mediciones y validaci\'on}
\label{sec:procedimiento_mediciones}

PLACEHOLDER: detallar el procedimiento utilizado para validar el
funcionamiento del sistema completo a partir de mediciones reales sobre
el canal HSR, en l\'inea con el objetivo espec\'ifico correspondiente
(Secci\'on~\ref{sec:objetivos_especificos}).
